\section{Problems}\label{sec:open}

Problems~\ref{prob:sixthree} and~\ref{prob:seventhree} together decide the
real problem at rank three, by Theorem~\ref{thm:rankthree}.

\begin{enumerate}

\item\label{prob:sixthree}
\emph{Prove Conjecture~\ref{conj:sixthree}: the stationarity system
\eqref{eq:statg}--\eqref{eq:statt} has no admissible solution at $(6,3)$ with
$0<\cv<1$.}

Such a proof closes $\Gtz(6,3)$ through Theorem~\ref{thm:rankthree}, which also
assumes that the cell has a minimizing counterexample
(Definition~\ref{def:extremalcex}).  Theorem~\ref{thm:variational} supplies that
assumption for the chart objective, on the two instances $\Gtz(5,3)$ and
$\Gtz(5,2)$ of Corollary~\ref{cor:corank}; by Remark~\ref{rem:objectives} a
minimizer of the chart objective need not minimize the design objective, and
nothing in the paper closes that gap.

The question is finite.  The active family $A^{+}$ covers the six
indices by \eqref{eq:coverage}, so it is one of the $2102$ families counted
in Proposition~\ref{prop:census}.  It has at least three members except in the
two-triple partition branch that Corollary~\ref{cor:e3} leaves undecided, and
setting that class aside leaves $2101$; it may be taken of cardinality
at most $15$ by Carath\'eodory's theorem \cite{Caratheodory1911}
(Proposition~\ref{prop:carath}), which leaves $2068$ (\S\ref{app:census}).  The
exclusions of \S\ref{sec:exclusions} settle $57$ classes of the full census of
$2136$, and the remainder is concentrated where they are silent.  For each class
the task is to show that \eqref{eq:statg}, \eqref{eq:statt}, the eigenvector
equations $S_{C_i}u_i=\cv u_i$, \eqref{eq:parseval} and the admissibility
disjunction \eqref{eq:argmax} have no common real solution with every
$t_c>0$, every $\lambda_i\ge0$ and $0<\cv<1$.

Any elimination must respect the system's multiplicities and its redundancy.  If
$\cv$ is an eigenvalue of $\SC$ of multiplicity $r$ at an active $C$, then the
multiplier attached to $C$ is a positive semidefinite matrix of trace one on
that eigenspace \eqref{eq:spectraplex}, and $C$ carries up to $r$ of the indices
of $A$ rather than one.  And \eqref{eq:coverage} is \eqref{eq:statg} contracted,
with $\mu=\cv$ following from it, so the constraints must not be counted as
independent.

\item\label{prob:seventhree}
\emph{Prove Conjecture~\ref{conj:seventhree}: the same at $(7,3)$.}

With Problem~\ref{prob:sixthree} this decides the real problem at rank three,
by Theorem~\ref{thm:rankthree}.  Enumeration is not available: the census has
$7\,011\,184$ classes, of which $5\,249\,381$ lie within the Carath\'eodory
bound $19$ (\S\ref{app:census}).  An argument uniform in the active family is
wanted, one using only the laws \eqref{eq:statt} and \eqref{eq:coverage}, the
positivity of the weights, and the combinatorics of a covering family, never
the family's identity.

The split tetrahedra of \S\ref{app:split} have value $\cv=1$ at this cell, and
their active family is the set of triples meeting three distinct classes, so its
cardinality is the third elementary symmetric function of the class sizes: $13$,
$17$ and $20$ at the profiles $(4,1,1,1)$, $(3,2,1,1)$ and $(2,2,2,1)$.  A
uniform argument must therefore exclude $\cv<1$ without excluding any of the
three.

\item\label{prob:sharpc}
\emph{Determine $\alpha_3(\C)$.}

Remark~\ref{rem:sharpc} gives
$\tfrac13\le\alpha_3(\C)\le3\bigl(1-\cos\tfrac{2\pi}{9}\bigr)=0.7018\ldots$,
the upper bound realized by a binding triple of the Hesse configuration of nine
equiangular lines in $\C^3$.  We conjecture that the upper bound is the truth.
At $k=2$ both ends meet, $\alpha_2(\C)=2-2/\sqrt3$; for $k\ge3$ the two ends of
\S\ref{sec:csharp} are all that is known.  Over $\R$ the same question is empty
below rank three and, if Problems~\ref{prob:sixthree}
and~\ref{prob:seventhree} go as conjectured, empty at rank three, so it first
has content at $\alpha_4(\R)$.

\item\label{prob:interlace}
\emph{Derandomize Theorem~\ref{thm:average}.}

Volume sampling \cite{DRVW2006} assigns $\pi_C=\det P[C]$ to a $k$-subset by
\eqref{eq:dpp}.  The measure is determinantal, carried by the bases of the
underlying matroid \cite{Lyons2003}, and it is strongly Rayleigh
\cite{BorceaBrandenLiggett2009}, so the mixture $q=\E_\pi[\chi_{\SC}]$ is
real-rooted \cite{AnariOveisGharan2015}.  Two statements would together prove
$\Gtz(m,k)$ at every rank at once:
\begin{enumerate}
\item[(a)] some $k$-subset $C$ has $\lmin(\SC)$ at least the least root of $q$;
\item[(b)] that least root is at least $1$.
\end{enumerate}

Neither is available.  Statement (a) does not follow from real-rootedness of
$q$ alone (Proposition~\ref{prop:rootedness}), and the interlacing method of
\cite{MSS2015} requires the whole convex family to be real-rooted, a property
of volume sampling that has not been verified.  Statement (b) must be exact
on the equality locus, by Proposition~\ref{prop:qone}: at a design where
every basis has $1$ in the spectrum of $\SC$ one has $q(1)=0$, with no
margin, and at the tetrahedron $q=(x-1)(x-4)^2$ (\S\ref{app:tetra}).

A route through this problem would be the first argument in the subject that
does not proceed cell by cell.  Nothing in this paper manages that.

\item\label{prob:structure}
\emph{Classify the extremal designs.}

Theorem~\ref{thm:corankone} classifies them at corank one: over each weight
vector there is one, given in \S\ref{app:corankone}, and \eqref{eq:tielaw} is a
complete description.  Beyond corank one the constructed families are the
split-class designs of \S\ref{app:split}, which have at most $k+1$ distinct
directions, and the splittings of the diamond of \S\ref{app:diamond}, which
has $k+2$.

Is that list complete at $(6,3)$ and $(7,3)$?  Is it finite, as a list of
families?  Does every extremal design have every basis active, that is
$\lmin(\SC)=1$ for every independent $k$-subset $C$?
Every known extremal design does: at the tetrahedron all four
triples are bases and all are active, and at the diamond the eight spanning
trees are bases and all eight have positive semidefinite singular gaps, the two
circuits being dependent and carrying no volume-sampling mass.  An affirmative
answer gives the hypothesis under which Proposition~\ref{prop:qone}
yields $q(1)=0$, and so is what Problem~\ref{prob:interlace} would have to
face at the equality locus.

\item\label{prob:higher}
\emph{Determine the true window in Theorem~\ref{thm:window}.}

Nothing in the proof of Lemma~\ref{lem:null}
refers to domination, so no lower bound on the window follows from it.  At
$k=2$ the
problem is decided at every size by Theorem~\ref{thm:ranktwo}, so the window is
vacuous there; at $k=3$ the bound is $7$ and the two cells $(6,3)$ and $(7,3)$
survive.  Is the window ever larger than $k+3$?  Equivalently: is there a rank
$k$ and a size $m>k+3$ at which $\Gtz(m,k)$ is strictly harder than at every
smaller size?  Corollary~\ref{cor:corank} settles corank at most two, so the
question is first posed at corank three.

\item\label{prob:effective}
\emph{Find the subset.}

The selection principle asserts that a subset exists; neither known route
exhibits one.  Volume sampling dominates only on average
(Theorem~\ref{thm:average}).  The maximal-volume principle
\cite{GT2001} does name a subset, and the guarantee it comes with is too
weak.  At uniform
weights that principle bounds a locally maximal-volume choice by
$\lVert A_C^{-1}\rVert^2\le1+k(m-k)$ for the matrix $A$ of rows $u_c=g_c/\sqrt m$,
which in the normalization of \eqref{eq:dominate} reads
\[
  \lmin(\SC)\;\ge\;\frac{m}{1+k(m-k)} .
\]
By \eqref{eq:deficit} the right side equals $1$ exactly when $k=1$ or $m=k+1$
and is strictly smaller otherwise, so the principle reproduces
Theorem~\ref{thm:corankone} and nothing beyond it.

Computing a maximal-volume subset is itself out of reach: the problem is
NP-hard and NP-hard to approximate within a constant factor
\cite{CivrilMagdonIsmail2009}, and the rank-revealing factorizations that
approximate it \cite{GuEisenstat1996} inherit the weaker guarantee.

The witnesses named below are Lean declarations, verified and covered by
the ledger of \S\ref{sec:axioms}.

At general weights the rule fails outright.  The volume of $C$ is $\pi_C$ of
\eqref{eq:dpp}; at a design at $(7,3)$ whose atom vectors are integral, whose
weights are integer multiples of $\tfrac1{60}$ and whose leverages are all at
least $2$, the maximum of $\pi_C$ over the thirty-five triples is attained at a
triple that does not dominate
(\texttt{Gtz.shadowDeterminant\_\allowbreak le\_\allowbreak shadowArgmaxSubset},
\texttt{Gtz.shadowArgmaxSubset\_\allowbreak not\_\allowbreak dominates}), while
another triple does
(\texttt{Gtz.shadowArgmaxDesign\_\allowbreak hasDominatingSubset}).  At that same design the
maximizer of the unweighted volume $\lvert\det M_C\rvert$ dominates strictly, so
what breaks the rule is the weighting by $\prod_{c\in C}t_c$ in \eqref{eq:dpp},
and Proposition~\ref{prop:rankfloor}, whose subset is the unweighted maximizer,
is untouched.

Whatever exhibits the subset must read the atom labels.  No triple dominates at
every design at $(6,3)$: two rational designs there each have a dominating
triple and share none
(\texttt{Gtz.no\_\allowbreak universal\_\allowbreak dominating\_\allowbreak subset},
\texttt{Gtz.exists\_\allowbreak designs\_\allowbreak with\_\allowbreak disjoint\_\allowbreak dominationSets}).  Invariant data
is not enough either: some design at $(6,3)$ is fixed atom-for-atom and
weight-for-weight by a nontrivial relabelling of its indices, has a dominating
triple, and has none that the relabelling fixes
(\texttt{Gtz.exists\_\allowbreak symmetry\_\allowbreak with\_\allowbreak no\_\allowbreak fixed\_\allowbreak dominatingSubset}), so every
rule computed from relabelling-invariant data fails there.  That witness is the
$(6,3)$ split tetrahedron of \S\ref{app:split} at $a=b=\tfrac18$, the
relabelling exchanging the two copies in each split direction, so the
obstruction sits on the equality locus.  Taking the atom of greatest leverage
fails at a design all of whose leverages exceed $1$
(\texttt{Gtz.heaviest\_\allowbreak atom\_\allowbreak can\_\allowbreak lie\_\allowbreak outside\_\allowbreak every\_\allowbreak dominatingSubset},
\texttt{Gtz.heavyPivotDesign\_\allowbreak allHeavy}).
Maximizing $\lmin(\SC)$ is not refuted and cannot be: its argmax dominates
whenever any triple does, and computing that argmax is the exhaustive search.

Local search is refuted at a bounded exchange radius.  Suppose a score has the
property that at every non-dominating $k$-subset some $k$-subset within $r$
swaps scores strictly higher; a finite argmax then proves $\Gtz(m,k)$.  At rank
three and $r\le2$ the property fails on all-heavy designs for $\lmin(\SC)$ at
both open cells, and at $(7,3)$ for the clipped trace, the trace of $\SC$ with
each eigenvalue replaced by its minimum with $1$, and for $\pi_C$
(\texttt{Gtz.not\_\allowbreak exchangeImprovesWithinRadiusHeavy\_\allowbreak of\_\allowbreak le\_\allowbreak two\_\allowbreak rankThree},
\texttt{Gtz.boundedRadiusExchangeHeavy\_\allowbreak refuted\_\allowbreak sevenThree}).  At $r=k$ the
radius constraint is vacuous.  The first two scores separate dominating from
failing subsets by a threshold, so at that radius their hypothesis is
$\Gtz(m,k)$ verbatim; for $\pi_C$ it is false with no radius bound at all.
Every witness above satisfies $\Gtz$ at its own cell.

Two questions therefore stand apart from the conjecture itself: does a subset of
maximal volume dominate at uniform weights, and, granting $\Gtz(m,k)$, is there
an algorithm polynomial in $m$ and $k$ that exhibits a dominating subset?

\end{enumerate}
